\section{Invariant-quotient monodromy proof}
\label{app:monodromy}

We give the linear-algebra argument behind
\cref{prop:quotient-identity,thm:monodromy-gap} in a form that includes
possible Jordan blocks at the unit multiplier.

Fix a periodic point \(z\) and abbreviate
\begin{equation}
 V=T_z\R^4,
 \qquad \alpha=dK_\epsilon(z),
 \qquad v=X_{K_\epsilon}(z),
 \qquad L=\Span\{v\},
 \qquad M=M_\epsilon.
 \label{eq:flag-notation}
\end{equation}
Since \(dK(X_K)=0\), \(L\subset\ker\alpha\).  Flow covariance at the full
return gives
\begin{equation}
 Mv=D\Phi^{T}(z)X_K(z)=X_K(\Phi^T(z))=v.
 \label{eq:flow-vector-fixed}
\end{equation}
Differentiated energy conservation gives
\begin{equation}
 \alpha\circ M=\alpha.
 \label{eq:covector-fixed}
\end{equation}
Therefore \(0\subset L\subset\ker\alpha\subset V\) is an invariant
filtration.  The induced map on \(L\) is the identity by
\cref{eq:flow-vector-fixed}; the induced map on
\(V/\ker\alpha\) is the identity by \cref{eq:covector-fixed}.

Variations tangent to the energy shell belong to \(\ker\alpha\).  Changing
the event time changes a variation by a multiple of \(v\).  Event
transversality therefore identifies the derivative of the two-dimensional
return with the induced map on \(\ker\alpha/L\).  Characteristic
polynomials multiply over invariant subspaces and quotients, so
\begin{equation}
 \chi_M(t)=(t-1)^2\chi_{D\Pi}(t).
 \label{eq:flag-factorization-app}
\end{equation}
No eigenbasis has been introduced.  The argument remains valid if the two
unit multipliers have a nontrivial Jordan block, or if their algebraic
multiplicity is four.

With the Hamiltonian sign convention used here, \(\ker\alpha=L^\omega\).
The symplectic form descends to a nondegenerate two-form on
\(L^\omega/L\).  Since \(M\) is symplectic, its induced map preserves this
form, and
\begin{equation}
 \det D\Pi=1.
 \label{eq:return-area}
\end{equation}
Taking traces in \cref{eq:flag-factorization-app} gives
\(\operatorname{tr}M=2+\operatorname{tr}D\Pi\).  Hence
\begin{equation}
 \det(I-D\Pi)
 =1-\operatorname{tr}D\Pi+\det D\Pi
 =4-\operatorname{tr}M.
 \label{eq:det-identity-app}
\end{equation}

The stored four-dimensional physical diagonal entries have flattened
indices \(0,7,14,21\); augmented parameter and period coordinates are not
included.  For each transcript the directed interval is
\begin{equation}
 [D_M]=
 \left[4-\sum_{j=0}^{3}M_{jj}^{+},
       4-\sum_{j=0}^{3}M_{jj}^{-}\right].
 \label{eq:directed-det-interval}
\end{equation}
The exact-fraction lower endpoints exceed three for all 202 jobs.  Their
18-place downward displays are those in \cref{eq:gap-minima}.  Maximum
interval widths are \(0.054493101512001146\) and
\(0.025036862429395394\) at 128 and 256 bits.  These widths suffice for the
strict inequality but do not satisfy the separate, still-open
Taylor-identity residual target.
